% p220_1; source PDF page 220; manually verified editable table.
\begin{tikzpicture}[x=1cm,y=-1cm,inner sep=0pt,outer sep=0pt]
\path[use as bounding box] (0,0) rectangle (16.000,3.200);
\draw[line width=.4pt] (0.000,0.000) rectangle (4.500,0.900);
\node[align=center,text width=4.240cm,font={\linespread{1}\fontsize{10}{13.50}\selectfont\bfseries}] at (2.250,0.450) {行为类型};
\draw[line width=.4pt] (4.500,0.000) rectangle (9.700,0.900);
\node[align=center,text width=4.940cm,font={\linespread{1}\fontsize{10}{13.50}\selectfont\bfseries}] at (7.100,0.450) {结论};
\draw[line width=.4pt] (9.700,0.000) rectangle (16.000,0.900);
\node[align=center,text width=6.040cm,font={\linespread{1}\fontsize{10}{13.50}\selectfont\bfseries}] at (12.850,0.450) {范例};
\draw[line width=.4pt] (0.000,0.900) rectangle (4.500,3.200);
\node[align=left,text width=4.240cm,font={\linespread{1}\fontsize{10}{13.50}\selectfont}] at (2.250,2.050) {欺骗或强迫他人自杀、自伤（形成支配力）};
\draw[line width=.4pt] (4.500,0.900) rectangle (9.700,3.200);
\node[align=left,text width=4.940cm,font={\linespread{1}\fontsize{10}{13.50}\selectfont}] at (7.100,2.050) {故意杀人罪、故意伤害罪的间接正犯};
\draw[line width=.4pt] (9.700,0.900) rectangle (16.000,3.200);
\node[align=left,text width=6.040cm,font={\linespread{1}\fontsize{10}{13.50}\selectfont}] at (12.850,2.050) {例1，皇帝赐大臣毒酒，赐死。\\例2，教唆、欺骗幼儿、精神病患者，使其自杀、自残};
\end{tikzpicture}
